\chapter*{Introduction générale}


Dans le cadre de mon Projet de Fin d’Études au sein de MEDIASSIVE, entreprise basée à Tanger et spécialisée dans le développement digital, les solutions web et l’innovation technologique, j’ai travaillé sur la conception, le développement et le déploiement d’une plateforme web nommée FinderrLink. Ce projet s’inscrit dans la vision de MEDIASSIVE, qui vise non seulement à proposer des services digitaux, mais aussi à lancer ses propres solutions SaaS internes répondant à des besoins réels du marché.

FinderrLink est une plateforme de mise en relation entre freelances, entreprises et agences. Elle permet de centraliser la gestion des profils professionnels, des missions, des candidatures, de la messagerie, des notifications et des tableaux de bord dans un seul environnement numérique. Elle vise ainsi à simplifier la recherche de talents, la publication des missions et le suivi des collaborations professionnelles.

Durant ce projet, j’ai occupé le rôle de chef d’équipe, ce qui m’a permis de participer aux aspects techniques et organisationnels. Mon intervention a porté sur l’analyse des besoins, la conception de l’architecture, le développement full-stack, la mise en place du système RBAC, la coordination des tâches et la préparation du déploiement. Le système RBAC permet de gérer les droits d’accès selon les rôles : freelance, entreprise, agence et administrateur \cite{owasp}.

Le développement de FinderrLink a présenté plusieurs défis, notamment la gestion de rôles différents, la structuration de la base de données, le suivi des candidatures, l’intégration de la messagerie, la mise en place des notifications et le déploiement dans un environnement serveur. Une particularité importante du projet réside dans le rôle hybride de l’agence, qui peut publier des missions comme une entreprise et postuler à des missions comme un freelance.

Ce projet m’a permis de renforcer mes compétences en développement web full-stack, en conception logicielle, en sécurité applicative, en gestion de base de données, en déploiement et en coordination d’équipe. Le présent rapport présente l’entreprise d’accueil, le contexte du projet, l’analyse des besoins, la conception, les technologies utilisées, la réalisation des modules, le déploiement, les tests et les perspectives d’amélioration de FinderrLink.

